\begin{itemize}
    \item Complementar el análisis exploratorio con indicadores financieros adicionales, con el fin de evaluar el rendimiento ajustado por riesgo de cada empresa.
    \item Incorporar variables fundamentales, tales como ingresos, utilidades, capitalización bursátil, indicadores macroeconómicos o tasas de interés, para comprender mejor los factores que explican el comportamiento de las acciones.
    \item Profundizar en el análisis de la relación entre precio y volumen mediante medidas estadísticas como correlaciones o modelos de regresión, permitiendo cuantificar el grado de asociación entre ambas variables.
    \item Realizar análisis por subperíodos o eventos relevantes (por ejemplo, pandemia de COVID-19, cambios en política monetaria o publicación de resultados financieros) para identificar cómo diferentes escenarios afectaron el comportamiento bursátil de las empresas.
    \item Utilizar los resultados obtenidos como punto de partida para construir modelos predictivos o estrategias de inversión que combinen rendimiento esperado, volatilidad y liquidez, permitiendo una toma de decisiones más informada y fundamentada.
\end{itemize}
